\documentclass[11pt]{article}

% Shared preamble for the SKJS soundness corpus (main paper + deep dives).
% Loaded via \input in each document's preamble. Documents may define
% additional packages, theorem environments, and macros AFTER this input,
% but must not redefine anything here.
\usepackage{amsmath}
\usepackage{amssymb}
\usepackage{amsthm}
\usepackage{mathtools}
\usepackage{stmaryrd}
\usepackage[margin=1in]{geometry}
\usepackage{tcolorbox}
\usepackage{booktabs}
\usepackage{array}
\usepackage{longtable}
\usepackage[hyphens]{url}
\urlstyle{tt}

% Long typewriter tokens (paths, identifiers) defeat hyphenation;
% allow inter-word stretch instead of overfull lines.
\emergencystretch=2em

\newtheorem{theorem}{Theorem}
\newtheorem{lemma}[theorem]{Lemma}
\newtheorem{definition}[theorem]{Definition}
\newtheorem{corollary}[theorem]{Corollary}
\newtheorem{axiom}[theorem]{Axiom}
\newtheorem{obligation}[theorem]{Obligation}

% Summary box for key definitions
\newtcolorbox{keydef}{
  colback=gray!5,
  colframe=gray!50,
  boxrule=0.5pt,
  arc=2pt,
  left=4pt,right=4pt,top=2pt,bottom=2pt
}

% Core semantic notation shared by every document
\newcommand{\sem}[1]{\llbracket #1 \rrbracket}
\newcommand{\Val}{\mathsf{Val}}

% Implementation anchor: file / class / function at the pinned revision
\newcommand{\impl}[1]{\par\noindent{\footnotesize\textsf{impl:} \texttt{#1}}}

% Single source for the modeled revision. Every document's \date uses this;
% the ledger and tooling read the same values.
\newcommand{\skjspinhash}{9108f339}
\newcommand{\skjspindate}{2026-07-02}
\newcommand{\skjspin}{Models revision \texttt{\skjspinhash} of \texttt{main} (\skjspindate)}



\newtheorem{finding}[theorem]{Finding}

% Core semantic notation (copied from the main paper)
\newcommand{\ObjVal}{\mathsf{ObjVal}}
\newcommand{\SType}{\mathsf{SType}}
\newcommand{\STypes}{\mathsf{STypes}}
\newcommand{\TType}{\mathsf{TType}}
\newcommand{\Clauses}{\mathsf{Clauses}}
\newcommand{\Clause}{\mathsf{Clause}}
\newcommand{\TInst}{\mathsf{TInst}}
\newcommand{\TInsts}{\mathsf{TInsts}}
\newcommand{\TParam}{\mathsf{TParam}}
\newcommand{\TFields}{\mathsf{TFields}}
\newcommand{\TObject}{\mathsf{TObject}}
\newcommand{\Tags}{\mathsf{Tags}}
\newcommand{\tagfn}{\mathsf{tag}}
\newcommand{\rtag}{\mathsf{rtag}}
\newcommand{\clause}[2]{\langle #1 \mid #2 \rangle}

% Operations under audit
\newcommand{\union}{\mathsf{union}}
\newcommand{\inter}{\mathsf{inter}}
\newcommand{\Env}{\Gamma}
\newcommand{\Payloadbool}{\{(\mathsf{bool},\mathsf{true}),
(\mathsf{bool},\mathsf{false})\}}

% Implementation anchors, as in the main paper
% Overrides the shared \impl with a raggedright variant (long anchors).
\renewcommand{\impl}[1]{\par\noindent{\footnotesize\raggedright
\textsf{impl:} \texttt{#1}\par}}
\setlength{\emergencystretch}{2em}

\title{Union and Intersection in SKJS:\\
an Audit of the Lattice Operations\\
{\large Deep Dive B for ``Refinement and Inference Growth in SKJS''}}
\author{}
\date{\skjspin}

\begin{document}

\maketitle

\begin{abstract}
The main soundness paper assumes that \texttt{TType::union}
over-approximates set union (\textbf{L-UNION}). It splits intersection
into non-bottom containment, conditional on faithful recursive bottoms
(\textbf{L-INTER}), and the bottom-implies-disjointness case
(\textbf{L-INTER-BOT}); together those cases say that
\texttt{TType::inter} over-approximates semantic intersection.
We check both operations at every layer of the type representation. Union
always contains the semantic union, and is exact at every layer observed by
the membership semantics; its cache is independent of the inference
substitution. Non-bottom intersection contains the semantic intersection in
the audited structural strata; stale \texttt{seenInter} reuse remains open.
The bottom clause fails in two witnessed cases. A failed function
signature intersection is treated as if the containing object branches were
disjoint, and distinct beyond-precision numeric spellings can denote the same
double while intersecting to $\bot$. In both cases the checker narrows a
reachable value to \texttt{never}. The remaining non-cache intersection cases
satisfy the required property. \texttt{seenInter} is known to retain results
computed while candidate growth is temporarily suppressed. The unresolved
question is whether reusing such a stale entry can cause an unsound acceptance
without also reusing a stale subtyping-cache entry
(\textbf{C-FREEZE-CACHE}).
\end{abstract}

\tableofcontents

%==========================================================================
\section{Purpose and Ledger Rows Addressed}
\label{sec:purpose}

This document is deep dive B in the series supporting the main paper
(\emph{Refinement and Inference Growth in SKJS}, same pinned revision).
Its lattice rows in the main paper's ``Assumptions and Proof
Obligations'' ledger are: \textbf{L-UNION} (union over-approximates set
union at all levels), established here \emph{exact at every stratum};
\textbf{L-INTER} (non-bottom containment, conditional on recursive
bottom-faithfulness); and
\textbf{L-INTER-BOT} (the $\bot$-implies-disjointness clause), which
is refuted in two strata. It also owns the cache obligation
C-FREEZE-CACHE, the remaining open part of L-INTER's implementation
coverage; its
proofs consume the canonical variable-denotation row TVAR-SEM (owner:
main paper).

\noindent
Note that L-INTER and L-INTER-BOT are the two cases of universal
intersection containment: if
$\inter(t_1,t_2) = \bot$ while $\sem{t_1}_\rho \cap \sem{t_2}_\rho \neq
\emptyset$, then the containment fails at that pair (nothing is in
$\sem{\bot} = \emptyset$), and conversely. The ``$\bot$ implies
disjointness'' phrasing is called out separately in the main paper only
because it is the form its refinement rules consume: \textsc{R-Eq}'s
true branch installs $\inter(\Env(x), \tau_e)$ directly, and
\textsc{R-EqField}, \textsc{R-Inst}, and \textsc{R-IsArray} \emph{drop an
object branch} (or the whole variable, for \textsc{R-Eq} on a
single-member type) exactly when an intersection comes back $\bot$.
Section~\ref{sec:findings} exhibits reachable programs where that drop is
wrong. The audit proceeds bottom-up through the representation strata of
the main paper's Section~3; the precise ledger edits are in
Section~\ref{sec:status}.
\impl{TTypeUnion.sk; TTypeInter.sk; TType.sk / STypes, Clause, Clauses,
TFields, PrimLiterals.sk}

%==========================================================================
\section{Recap: Semantics and the Shape of the Two Operations}
\label{sec:recap}

We reuse the main paper's runtime model and denotations without change:
values carry one of nine tags ($\rtag$); simple types are units, literal or
full primitives, and object types with branch sets; $\sem{\cdot}$ is the
literal-closure semantics for primitives and the extensible, get-only
membership semantics for object branches (main paper, Section
``Internal Type Representation''). A $\TType$ is a pair $(I, C)$ of an instance layer
$I \in \TInsts$ and clauses $C \in \Clauses$ in DNF, with
$\sem{(I,C)}_\rho = \sem{I}_\rho \cup \sem{C}_\rho$; a clause
$\clause{V}{S}$ denotes $\bigcap_{v\in V}\sem{v}_\rho \cap \sem{S}_\rho$
with $S \in \STypes \cup \{\top\}$, and $\STypes$ holds at most one
$\SType$ per tag. All denotations are indexed by the substitution state
$\rho$; we suppress $\rho$ when an argument is pointwise in it.

\paragraph{Union's shape.}
\texttt{TType::union} short-circuits top/bottom, then works
\emph{structurally}: instance sets are set-unioned unexpanded, and clauses
are merged by var-set key, colliding clauses merging their $\STypes$
per tag (literal-set union for primitives, branch-set union for objects).
Crucially, no arm of this computation expands an instance or reads a
variable binding: union is a pure function of its operands' syntax.
\impl{TTypeUnion.sk / TType::union\_, Clauses::union, STypes::union,
TObject::union, TInsts::union}

\paragraph{Intersection's shape.}
\texttt{TType::inter} short-circuits top/bottom, then handles bare
variables specially (constraint attachment or an unexpanded
\texttt{TInter}), then \emph{expands both operands to clauses} under the
current state (sealing variables it reads covariantly), runs two
speculative boxed-primitive subtype checks, and finally distributes:
pairwise clause intersection, per-tag $\STypes$ intersection, pairwise
object-branch intersection, field/index/call/new/brand intersection inside
a branch pair. Bottom can be produced at four depths: literal sets, the
per-tag map, a branch pair, and the whole clause product.
\impl{TTypeInter.sk / TType::inter\_, Clauses::inter, STypes::inter,
TObject::inter, TFields::inter}

\paragraph{Method.}
For each stratum we state the containment against the semantics of the
main paper, prove it or reduce it to a named assumption, and separately
classify every $\bot$-producing site as \emph{semantic} (disjointness
holds) or \emph{heuristic} (it does not). Claims about observable behavior
were tested against the pinned checker binary
(\texttt{target/host/dev/skjs}, run via
\texttt{make check-file-in-docker}); the probe protocol and pointers to
the executed witnesses are Section~\ref{sec:findings}.

%==========================================================================
\section{Union, Level by Level}
\label{sec:union}

\subsection{Literal sets and primitive members}

\begin{lemma}[Literal-Set Union is Exact]
\label{lem:lit-union}
For literal sets $L_1, L_2$ of the same kind $k$:
$\sem{\mathsf{prim}(k, L_1)} \cup \sem{\mathsf{prim}(k, L_2)} =
 \sem{\mathsf{prim}(k, L_1 \cup L_2)}$,
where $L_1 \cup L_2$ is the implementation's union of spelling sets.
\end{lemma}
\begin{proof}
By the closure semantics (main paper, Definition of literal denotation),
$\sem{\mathsf{prim}(k,L)} = \bigcup_{\ell \in L} \sem{\mathsf{prim}(k,
\{\ell\})}$; union of spelling sets commutes with the outer union. No
normalization or injectivity assumption is needed: even $\approx$-equivalent
duplicate spellings only make the two sides equal unions of equal parts.
\impl{PrimLiterals.sk / union (keyed set union of spellings)}
\end{proof}

\begin{lemma}[TPrimValue and TPrimUnit Union are Exact]
\label{lem:prim-union}
\texttt{TPrimValue::union} on same-kind operands returns: the full
primitive if either side is full ($L = \top$ absorbs); otherwise the
literal-set union, except that a boolean literal set reaching two elements
widens to full \texttt{boolean}. Each case is semantically exact.
Same-tag units are identical, so \texttt{STypes::union}'s unit arm
(return the left operand) is exact.
\end{lemma}
\begin{proof}
Full absorbs because $\sem{\mathsf{prim}(k,\top)} \supseteq
\sem{\mathsf{prim}(k,L)}$ for any $L$ (every literal denotes same-tag
values). The literal case is Lemma~\ref{lem:lit-union}. For boolean: the
only spellings of kind \texttt{bool} are \texttt{true} and \texttt{false},
so a two-element set already denotes all of
$\Payloadbool$; the widening is an equality. Units: a tag determines its
unit, and \texttt{mapUnion} pairs members of identical tags only.
\impl{TTypeUnion.sk / TPrimValue::union; STypes::union (TPrimUnit arm)}
\end{proof}

\subsection{Object branches: set union}

\begin{lemma}[Object Union is Exact]
\label{lem:obj-union}
$\sem{\mathsf{object}(B_1)} \cup \sem{\mathsf{object}(B_2)} =
\sem{\mathsf{object}(B_1 \cup B_2)}$.
\end{lemma}
\begin{proof}
Membership in $\mathsf{object}(B)$ is an existential over branches (main
paper, Definition of object-type semantics), and existentials distribute
over set union of the branch set. \texttt{TObject::union} is exactly
$B_1 \cup B_2$; the \texttt{unexpandedType} hint is reset to \texttt{None},
which is irrelevant to $\sem{\cdot}$ (the main paper excludes it from
semantics). No branch is merged, so no precision or soundness question
arises at this level.
\impl{TTypeUnion.sk / TObject::union (branches.union)}
\end{proof}

The nontrivial object-\emph{merging} code (\texttt{TField::union} and
the per-read merge loops that apply it) is \emph{not} on the
$\union(t_1,t_2)$ path: it is a separate client family used to resolve
a read against a multi-branch object. We audit it in
Section~\ref{sec:tfields-union} because the main paper's
\textsc{R-EqField} proof leans on the resolved read type it produces,
and L-UNION is stated ``at all levels''.

\subsection{STypes, Clause, Clauses, TInsts}

\begin{lemma}[STypes Union is Exact]
\label{lem:stypes-union}
$\sem{S_1} \cup \sem{S_2} = \sem{\STypes\mathord{::}\union(S_1,S_2)}$.
\end{lemma}
\begin{proof}
\texttt{mapUnion} keeps members whose tag appears on one side and merges
same-tag pairs with Lemmas~\ref{lem:prim-union}
and~\ref{lem:obj-union}. Since $\sem{S} = \bigcup_{s \in S}\sem{s}$ and
tags partition $\Val$ (main paper, Lemma A0/A1), grouping the union by tag
is exact, and each per-tag merge is exact.
\impl{TTypeUnion.sk / STypes::union; InternalKeyedArrayMap.sk / mapUnion}
\end{proof}

\begin{lemma}[Clause Collision Merge is Exact]
\label{lem:clause-union}
For clauses with equal var-sets $V$:
$\sem{\clause{V}{S_1}} \cup \sem{\clause{V}{S_2}} =
 \sem{\Clause\mathord{::}\union}$, where the merge yields
$\clause{V}{\top}$ if either operand's $\STypes$ is $\top$ and
$\clause{V}{\STypes\mathord{::}\union(S_1,S_2)}$ otherwise.
\end{lemma}
\begin{proof}
Writing $A = \bigcap_{v\in V}\sem{v}$, the left side is
$(A \cap \sem{S_1}) \cup (A \cap \sem{S_2}) = A \cap (\sem{S_1} \cup
\sem{S_2})$ by distributivity of $\cap$ over $\cup$; then
Lemma~\ref{lem:stypes-union}, or $\sem{S_i} = \Val$ absorbing when a side
is $\top$. This is the ``key-collision widening'' of the operand pair: it
is not a widening at all semantically — the collision merge is exact.
\impl{TTypeUnion.sk / Clause::union}
\end{proof}

\begin{lemma}[Clauses Union is Exact; Top-Collapse is Sound]
\label{lem:clauses-union}
$\sem{C_1} \cup \sem{C_2} = \sem{\Clauses\mathord{::}\union(C_1,C_2)}$,
and the constructor invariant (a top clause, if present, is alone) is
preserved.
\end{lemma}
\begin{proof}
$\sem{C} = \bigcup_{c \in C}\sem{c}$; \texttt{mapUnion} keyed by var-set
keeps unmatched clauses and merges matched ones by
Lemma~\ref{lem:clause-union}, so the result denotes the union of all
clause denotations on both sides. \texttt{Clauses::create} then collapses
to $\{\clause{\emptyset}{\top}\}$ if any clause is top; this replaces the
set with a superset ($\sem{\clause{\emptyset}{\top}} = \Val$), preserving
the containment (and losing exactness only in the case where the union
already contained $\Val$, so not even there).
\impl{TTypeUnion.sk / Clauses::union; TType.sk / Clauses::create}
\end{proof}

\begin{lemma}[Instance-Layer Union is Exact]
\label{lem:insts-union}
$\sem{I_1}_\rho \cup \sem{I_2}_\rho = \sem{I_1 \cup I_2}_\rho$ at every
state $\rho$.
\end{lemma}
\begin{proof}
$\sem{I}_\rho = \bigcup_{i \in I}\sem{i}_\rho$ and \texttt{TInsts::union}
is set union of unexpanded instances; no instance is rewritten, so the
identity holds pointwise in $\rho$. This is the design point the
implementation comments call out: union keeps variables and named
instances side by side rather than resolving them.
\impl{TTypeUnion.sk / TInsts::union}
\end{proof}

\subsection{The union theorem}

\begin{theorem}[L-UNION, Clause and Instance Strata]
\label{thm:union}
For all $t_1, t_2$ and every state $\rho$:
$\sem{t_1}_\rho \cup \sem{t_2}_\rho \subseteq
 \sem{\union(t_1,t_2)}_\rho$.
Moreover the containment is an equality except possibly at the top/bottom
short-circuits (where it is an equality too: $t \vee \top = \top$,
$t \vee \bot = t$) — i.e.\ union is semantically exact for the membership
semantics.
\end{theorem}
\begin{proof}
If either operand is top or bottom, the short-circuit returns the exact
answer. Otherwise
$\union(t_1,t_2) = (I_1 \cup I_2,\ \Clauses\mathord{::}\union(C_1,C_2))$
and
$\sem{(I,C)} = \sem{I} \cup \sem{C}$ splits the claim into
Lemmas~\ref{lem:insts-union} and~\ref{lem:clauses-union}. The equality
cache check \texttt{ty1 == ty2} returns \texttt{ty1}, exact. Since every
step is either a pure set operation or one of the exact merges above, and
none reads $\rho$, the result holds at every state uniformly — a fact we
reuse for the cache in Section~\ref{sec:caches}.
\impl{TTypeUnion.sk / TType::union, union\_}
\end{proof}

Two remarks. First, exactness is relative to the main paper's
\emph{membership} semantics, in which function signatures are opaque;
the signature bodies that \texttt{TFun::union} builds (parameters
intersected contravariantly, \texttt{this} parameters intersected
through the \texttt{ThisParam} cases — attach-only absorbing, an
undeclared \texttt{this} acting as the intersection identity
(\texttt{TTypeUnion.sk:286-307}) — returns unioned) never enter
$\sem{\cdot}$ and are part of the subtyping surface instead — and
\texttt{TFun::union} is not reachable from $\union(t_1,t_2)$ anyway
(only from the per-read merging clients below). Second, \texttt{unionN}
of an empty sequence returns $\bot$, the exact empty union.

\subsection{The branch-merging clients: per-read merges and the
callability join}
\label{sec:tfields-union}

A multi-branch object is never collapsed into a summary branch; instead
each \emph{read} merges what the branches offer for that one lookup,
and every read path requires the key to be served by \emph{all}
branches (a lookup any branch cannot serve returns \texttt{None} and is
rejected upstream). The paths: \texttt{TObject::maybeGetType} resolves
the field per branch and unions the resolved types
(\texttt{TType.sk:1008-1022}); \texttt{resolveLiteralAccess} and
\texttt{maybeGetSymbol} merge the per-branch \texttt{TField}s with
\texttt{TField::union} (\texttt{TType.sk:1259-1279, 1115-1134});
accessible index lookups fold \texttt{TField::unionIndex}
(\texttt{TType.sk:1288-1305}); and the call/new surface is the
\emph{join} \texttt{TObject::modeOf} over the \texttt{Callability}
lattice with \texttt{TFuns::union} on signatures
(\texttt{TType.sk:1136-1204}). These feed the resolved read types of
\textsc{R-EqField}. The correctness statements:

\begin{lemma}[Per-Read Merges Over-Approximate]
\label{lem:tfields-union}
Let $o \in \sem{\mathsf{object}(B)}$, say $o \in \sem{b_1}$ with
$b_1 \in B$.
(1) If every branch resolves field $f$, then $o$'s value at $f$ (when
present) lies in the unioned resolved type returned by
\texttt{maybeGetType}, and $f$ is required of $o$ only if $b_1$
requires it.
(2) The same holds for the merged \texttt{TField} of
\texttt{resolveLiteralAccess}/\texttt{maybeGetSymbol}
(\texttt{TField::union}: get-union, $\mathsf{opt}_1 \vee
\mathsf{opt}_2$) and for folded index gets
(\texttt{TField::unionIndex}).
(3) \texttt{modeOf} returns the lattice join of the branch modes, so
$o$ satisfies the returned mode; in particular
$\mathsf{CPresent}$ is returned only when every branch is
$\mathsf{CPresent}$, and \texttt{maybeGetCall}/\texttt{maybeGetNew}
accept only that case.
\end{lemma}
\begin{proof}
(1) $b_1$'s resolved type for $f$ is one disjunct of the returned
\texttt{unionN}, and $o$'s value at $f$ lies in it ($o \in \sem{b_1}$,
field or covering-index clause of the object semantics); containment by
Theorem~\ref{thm:union}. Requiredness: the union is over per-branch
resolutions, each carrying its own optionality into the merged
optionality disjunction, so a key optional in $b_1$ stays optional.
(2) \texttt{TField::union} unions gets and or-s optionality — the
constraint of the merged field on $o$ is weaker than $b_1$'s
(get-containment by Theorem~\ref{thm:union}; a required merged field
is required on both sides, in particular on $b_1$).
(3) By case analysis on \texttt{modeOf}'s fold
(\texttt{TType.sk:1162-1204}): any $\mathsf{CUnknown}$ branch, or
mixed $\mathsf{CAbsent}$/positive branches, join to
$\mathsf{CUnknown}$ (no claim); all-$\mathsf{CAbsent}$ joins to
$\mathsf{CAbsent}$, which $o$ satisfies from $b_1$'s side; a positive
join ($\mathsf{CFunction}$, or $\mathsf{CPresent}$ with the
\texttt{TFuns::union} of the branch signatures) asserts
$[[\mathrm{Call}]]$, which $b_1$ (positive, since no
$\mathsf{CAbsent}$/$\mathsf{CUnknown}$ was seen) already asserted of
$o$; signatures are opaque to membership. The partial signature union
\texttt{modeOf} also returns is consumed only by recovery diagnostics,
never as a sound signature (the code comment states this;
\texttt{maybeGetCall} requires $\mathsf{CPresent}$,
\texttt{TType.sk:1206-1214}).
\emph{Set-types and access.} Membership never reads $\mathsf{set}$
(get-only semantics) — \texttt{TField::union} intersects set-types
(deferred as an unexpanded intersection under the
\texttt{tenv.expandInter} mode flag), the correct dual for
write-through-the-merge, but nothing here depends on it; and
\texttt{FieldAccess::combine} is most-restrictive, which only
\emph{restricts} programs (main paper, Definition of object semantics,
Never/Phantom discussion).
\impl{TType.sk / TObject::maybeGetType, resolveLiteralAccess,
maybeGetSymbol, modeOf, callMode, newMode, maybeGetCall, maybeGetNew,
maybeGetAccessibleIndexForKeyKind;
TTypeUnion.sk / TField::union, unionIndex, TFuns::union}
\end{proof}

Unlike the branch-set union of Lemma~\ref{lem:obj-union}, the per-read
merge is a genuine over-approximation (field gets are unioned across
branches), which is why the implementation uses it only for reads,
never as the representation of a union type.

%==========================================================================
\section{Intersection, Level by Level}
\label{sec:inter}

We now audit \texttt{TType::inter} in the same bottom-up order. The
summary: the primitive stratum is exact \emph{modulo spelling injectivity}
(\textbf{L-LIT}(ii) — proved for decoded strings and within-precision
numbers, violated for beyond-precision numeric spellings); the per-tag
assembly in \texttt{STypes::inter} is exact; inside an object branch
pair, fields, indexes, the brand, and the callability meet outside its
signature arm are sound (the brand conflicts, the array-branded
callability conflicts, and the no-shared-tag cases being the
\emph{semantic} bottoms), but the signature arm of the callability meet
converts a failed signature intersection into a branch drop, which is
heuristic. The instance layer is sound. Throughout,
\texttt{inter\_} expands its operands as it runs (covariant reads seal
variables), an evolution $\rho \to \rho'$; the operands transport across
it by \textbf{L-EXP} (expansion preserves or enlarges semantics; imported
from deep dive C), and the containments below, being pointwise in the
state, hold at $\rho'$.

\subsection{Literal sets and primitive members}
\label{sec:inter-prim}

\begin{lemma}[Literal-Set Intersection is Exact Modulo L-LIT(ii)]
\label{lem:lit-inter}
Assume distinct spellings in $L_1 \cup L_2$ denote
$\approx$-inequivalent values (\textbf{L-LIT}(ii)). Then
$\sem{\mathsf{prim}(k,L_1)} \cap \sem{\mathsf{prim}(k,L_2)} =
 \sem{\mathsf{prim}(k, L_1 \cap L_2)}$, with $\bot$ (\texttt{None}) when
the spelling intersection is empty; under the assumption that bottom is
semantic. Without the assumption, both the containment and the
disjointness reading fail (Finding~\ref{find:spelling}).
\end{lemma}
\begin{proof}
($\supseteq$ direction of the equality is immediate from
$L_1 \cap L_2 \subseteq L_i$.) For $\subseteq$: let
$v \approx \delta_k(\ell_1) \approx \delta_k(\ell_2)$ with
$\ell_i \in L_i$ ($v$ is in its own $\approx$-domain since literal
denotations exclude NaN). By transitivity
$\delta_k(\ell_1) \approx \delta_k(\ell_2)$, so by injectivity
$\ell_1 = \ell_2 \in L_1 \cap L_2$. The implementation intersects
\emph{spelling} sets, so this is exactly its computation. If injectivity
fails — as for beyond-precision numeric spellings, which alias one
double while remaining distinct lexemes — two spellings of the same
value make $L_1 \cap L_2 = \emptyset$ while the semantic intersection is
inhabited. For the string kind injectivity holds outright: escape
sequences are decoded at parse time
(\texttt{JSStringEscape.sk / decodeJSStringLiteral}, called from the
literal-building sites in \texttt{ParserExpr.sk} and
\texttt{ParserUtils.sk}), so the stored lexeme \emph{is} the denoted
string and spelling intersection coincides with value intersection.
\impl{PrimLiterals.sk / inter; TTypeInter.sk / TPrimValue::inter;
JSStringEscape.sk / decodeJSStringLiteral}
\end{proof}

\begin{lemma}[TPrimValue/TPrimUnit Intersection]
\label{lem:prim-inter}
On the \texttt{STypes::inter} path operands always have equal kinds (the
map is tag-keyed), and: full absorbs ($\top \cap L = L$, exact); two
literal sets reduce to Lemma~\ref{lem:lit-inter}; same-tag units are
equal, and returning the left one is exact. The kind-mismatch
\texttt{None} arm is unreachable from \texttt{STypes::inter} and, where
reachable (the \texttt{overlaps} client), is semantic by tag
disjointness (A1).
\impl{TTypeInter.sk / TPrimValue::inter, STypes::inter}
\end{lemma}

\subsection{The per-tag assembly: STypes::inter}
\label{sec:inter-stypes}

The correct per-tag decomposition of an intersection of unions is
\[
\sem{S_1} \cap \sem{S_2}
= \bigcup_{k \in \Tags}\ \sem{S_1[k]} \cap \sem{S_2[k]},
\]
by tag disjointness (A1): tags present on only one side contribute
nothing, and each shared tag contributes its per-tag intersection. The
sound and exact assembly is therefore: intersect per shared tag,
\emph{drop} the tags whose per-tag intersection is empty, and return
$\bot$ iff \emph{no} tag survives.

The implementation realises exactly this. \texttt{STypes::inter} builds
the result with \texttt{mapInterOption}, whose combiner returns
\texttt{Some(f(x,y))} on a shared tag: an empty per-tag intersection
(\texttt{f} returns \texttt{None}) is \emph{dropped} from the map, not
treated as the zero of the merge, and unmatched tags are dropped likewise;
the wrapping \texttt{if (!stypes.isEmpty())} guard then returns $\bot$
exactly when no tag survives. Hence
$\STypes\mathord{::}\inter(S_1,S_2) = \bot$ iff every shared tag has an
empty per-tag intersection, i.e.\ iff $\sem{S_1}\cap\sem{S_2}=\emptyset$:
the assembly is exact and its bottoms are semantic. For example
$\{\mathsf{str}{:}\texttt{"a"},\ \mathsf{num}{:}\{0\}\} \cap
\{\mathsf{str}{:}\texttt{"a"},\ \mathsf{num}{:}\{1\}\}$ drops the empty
\texttt{num} tag and keeps \texttt{str}${:}\texttt{"a"}$, so
\texttt{("a" | 0) \& ("a" | 1) = "a"}, not \texttt{never}.
\impl{TTypeInter.sk / STypes::inter; InternalKeyedArrayMap.sk /
mapInterOption}

With the assembly exact, the stratum's soundness reduces to the per-tag
lemmas (\ref{lem:prim-inter} for primitives,
Section~\ref{sec:inter-object} for objects).

\subsection{Object branch pairs}
\label{sec:inter-object}

\texttt{TObject::inter} forms all pairwise branch intersections and
returns $\bot$ iff \texttt{TFields::inter} is \texttt{None} on every
pair.
This outer shape is exactly right:

\begin{lemma}[Pairwise Branch Product]
\label{lem:obj-inter}
If for all $b_1 \in B_1, b_2 \in B_2$ the branch-pair operation
satisfies $\sem{b_1} \cap \sem{b_2} \subseteq \sem{\inter(b_1,b_2)}$
(reading $\sem{\texttt{None}} = \emptyset$), then
$\sem{\mathsf{object}(B_1)} \cap \sem{\mathsf{object}(B_2)} \subseteq
 \sem{\TObject\mathord{::}\inter}$, and the $\bot$ result is semantic
iff every pair's $\bot$ is.
\end{lemma}
\begin{proof}
$o$ in both objects picks $b_1, b_2$ with $o \in \sem{b_1} \cap
\sem{b_2}$; the surviving pair (or the contradiction with its dropped
status) is immediate, since the result's branch set contains
$\inter(b_1,b_2)$ whenever it is not \texttt{None}.
\impl{TTypeInter.sk / TObject::inter}
\end{proof}

One guard sits on this product: when the pairwise crossing of two
multi-branch objects exceeds the representation cap
(\texttt{kMaxObjectBranches}), \texttt{TObject::inter} raises a fatal
diagnostic (TS2590, ``too complex to represent'';
\texttt{TTypeInter.sk:609-624}) instead of materializing the
exponential branch set. That is a \emph{rejection}, not a $\bot$
result — conservative for soundness, a completeness bound only — so it
does not enter the bottom inventory.

So the object stratum reduces entirely to one branch pair, audited
constraint by constraint against the main paper's object semantics. Fix
$o \in \sem{b_1} \cap \sem{b_2}$ and let $b$ be the merged branch.

\paragraph{Brand (sound; semantic bottom).}
Conflicting brands $\mathsf{Some}(\mathsf{true})$ vs
$\mathsf{Some}(\mathsf{false})$ return \texttt{None}: semantic, since
$\mathsf{isArrayV}(o)$ cannot be both. Otherwise the known side (if any)
is kept, which $o$ satisfies. The same argument covers the refinement
helper \texttt{TFields::interIsArray} used by \textsc{R-IsArray}: setting
an unknown brand constrains $o$ only per the predicate already
established, and its mismatch bottom is semantic.
\impl{TTypeInter.sk / TFields::inter (isArray arms), interIsArray}

\paragraph{Named fields (sound).}
\label{sec:inter-fields}
The field maps merge with: both sides present $\to$
\texttt{TField::inter} (get intersected — or deferred as an unexpanded
intersection when \texttt{expandInter} is off, which denotes the same
intersection by the \texttt{TInter} interpretation — set unioned,
$\mathsf{opt}_1 \wedge \mathsf{opt}_2$, access combined); present on one
side $\to$ kept, intersected with the other side's covering index if one
exists (\texttt{interIndex}: get intersected, optionality preserved since
indexes are optional). Soundness for $o$: a key required on either side
is present ($o$ satisfies that side), so
$\mathsf{opt}_1 \wedge \mathsf{opt}_2$ is respected; when the key is
present, its value lies in each side's applicable get-constraint — the
named field's get on the side that names it, the index's get on a side
that covers it without naming it (the index clause of the object
semantics applies precisely because the key is \emph{not named there}) —
hence in the intersection of the collected gets, recursively by
L-INTER/L-INTER-BOT
at the field-type level. Access combination is most-restrictive
(\texttt{FAPhantom} absorbs into phantom, distinct-owner privates give
\texttt{FANever}); phantoms weaken the value constraint to
presence-only, a widening, and \texttt{FANever} remains an ordinary
membership constraint — both restrict-only at access sites, per the main
paper. Symbol-keyed fields are the same argument against the symbol
index.
\impl{TTypeInter.sk / TFields::interNamedFields, TField::inter,
TField::interIndex; TType.sk / FieldAccess::combine}

\paragraph{The set direction (membership-irrelevant; flagged).}
\texttt{TField::inter} \emph{unions} the set-types. For the get-only
membership semantics this is vacuous — $\sem{b}$ never reads
$\mathsf{set}$ — so L-INTER is unaffected; that is the whole proof. The
choice is the one that makes the merged field a subtype of both operands
under setter contravariance (the merged branch must accept every write
each operand accepts). But as a \emph{write-through} interface it is the
permissive direction: a write accepted because only $b_1$'s set-type
allows it can break $b_2$'s get-assumption on the same key. That is
a field-write concern --- membership is snapshot (main paper, object
semantics), writes are governed by subtyping --- so we do not count it
against L-INTER; deep dive H audits this direction against the
write-authorization semantics.
\impl{TTypeInter.sk / TField::inter (set => TType::union)}

\paragraph{Indexes (sound, via the covering semantics).}
Per kind, one side absent keeps the other (an absent index is the
unconstrained index, and intersecting constraints keeps the stronger
side); both present intersect gets. Then the number index is replaced by
$\mathsf{idx}_{\mathsf{num}} \sqcap \mathsf{idx}_{\mathsf{str}}$ when
both survive. Soundness for $o$: take a prop $p$ of $o$ covered by a
merged index and not named in $b$. Since $b$'s names are the union of
the sides' names, $p$ is not named on either side, so each side's
covering index (string indexes covering numeric-kind keys as in the main
paper's covering semantics) constrains $\mathsf{get}(o,p)$; the merged
get is the intersection of exactly those constraints — including the
$N \sqcap S$ adjustment, since a numeric-kind key not named anywhere is
covered by \emph{both} the effective number index and the string index
of each side that has them. No index arm returns bottom.
\impl{TTypeInter.sk / TIndexes::inter, TFields::interOptIndex}

\paragraph{Call/new: the callability meet.}
The call and new members are values of the four-point
\texttt{Callability} lattice
($\mathsf{CPresent}(\mathit{sig}) \le \mathsf{CFunction} \le
\mathsf{CUnknown}$, $\mathsf{CAbsent} \le \mathsf{CUnknown}$;
\texttt{TType.sk:2189-2237}), read by membership semantics as:
$\mathsf{CAbsent}$ — no member implements $[[\mathrm{Call}]]$
(resp.\ $[[\mathrm{Construct}]]$); $\mathsf{CFunction}$ and
$\mathsf{CPresent}$ — every member does; $\mathsf{CUnknown}$ — no
constraint; signatures opaque. Both arms of \texttt{TFields::inter} meet
the modes with \texttt{TFields::interCallability}
(\texttt{TTypeInter.sk:474-502}), arm by arm:
$(\mathsf{CUnknown}, m) \to m$ is exact (one side constrains nothing);
$(\mathsf{CAbsent}, \mathsf{CAbsent}) \to \mathsf{CAbsent}$ and
$(\mathsf{CFunction}, \mathsf{CFunction}) \to \mathsf{CFunction}$ are
exact; $(\mathsf{CFunction}, \mathsf{CPresent}) \to \mathsf{CPresent}$
is exact (both sides assert $[[\mathrm{Call}]]$; the signature is the
more precise claim). The mixed arm
$(\mathsf{CAbsent}, \mathsf{CPresent}/\mathsf{CFunction})$ is
brand-aware: off the array brand it \emph{keeps the positive mode} — a
sound over-approximation of a contradiction (a member cannot both
implement and not implement $[[\mathrm{Call}]]$), deliberately not
$\bot$ because $\mathsf{CAbsent}$ on nominal instances is a brand-like
approximation that structural width subtyping can falsify (the two-step
callable-to-class conversion documented in the code at
\texttt{isSubCallability}), so a
$\bot$ here would collapse inhabited hybrid types like
\texttt{(() => void) \& C}. Under the array brand
($\mathsf{isArray} = \mathsf{Some(true)}$) the same arm returns
\texttt{None} (\texttt{TTypeInter.sk:485-492}): \emph{semantic},
because an array-branded member is an array exotic object, which
implements neither $[[\mathrm{Call}]]$ nor $[[\mathrm{Construct}]]$
(ECMA-262 \S10.4.2), contradicting the positive mode — conditional on
the array brand's membership invariant, which subtyping leaves
unvalidated (registry row F-BRAND); a value accepted into an
array-branded view would make this bottom spurious
through that row, not through this arm. The meet also preserves the
constructible-implies-$[[\mathrm{Call}]]$ construction invariant
pairwise (\texttt{TType.sk:2273-2281}): it never invents a
constructible mode from $\mathsf{CUnknown}$ inputs. The one unsound
arm is $(\mathsf{CPresent}, \mathsf{CPresent})$: it delegates to
\texttt{TFuns::inter}, which intersects the signatures by trying
\texttt{TFun::isSub} in both directions and returning the smaller side,
and when \emph{neither} signature subsumes the other it returns
\texttt{None} — which \texttt{interCallability} propagates and
\texttt{TFields::inter} converts into dropping the whole branch pair.
``Neither is a subtype of the other'' is not disjointness: under the
membership semantics signatures are opaque, and even behaviorally, one
function value can satisfy two incomparable signatures (any overload;
e.g.\ a \texttt{(x: unknown) => void} value inhabits both
\texttt{(s: string) => void} and \texttt{(n: number) => void}).

\begin{finding}[Signature Bottom: L-INTER-BOT is Violated at the
Signature Arm of the Callability Meet]
\label{find:tfun}
$\inter$ of two single-branch function types with incomparable
signatures is $\bot$ although their semantic intersection is inhabited.
Empirically confirmed (probe~2, Section~\ref{sec:findings}): in
\texttt{if (k === g)} with \texttt{k: (s: string) => void} and
\texttt{g: (n: number) => void}, the checker narrows \texttt{k} to
\texttt{never} and accepts \texttt{const z: never = k}; no
no-overlap warning is emitted (the \texttt{overlaps} helper itself
correctly says reference types always overlap, in direct tension with
\texttt{TFun::inter}). The same conversion makes the type-level
\texttt{F = ((s: string) => void) \& ((n: number) => void)} expand to
\texttt{never} (probe~4), so nothing is assignable to \texttt{F}
(spurious rejection — the same defect surfacing as incompleteness).
The sound repair is to keep the branch with an overload-like or
widened signature (even $\mathsf{CPresent}$ with either operand's
signature is membership-sound, signatures being opaque), never to drop
it.
\impl{TTypeInter.sk / TFun::inter, TFuns::inter, TFields::interCallability,
TFields::inter (call/new arms); SType::overlaps (contrast)}
\end{finding}

\subsection{Clause and Clauses intersection}

\begin{lemma}[Clause Pair]
\label{lem:clause-inter}
$\sem{\clause{V_1}{S_1}} \cap \sem{\clause{V_2}{S_2}} \subseteq
\sem{\clause{V_1 \cup V_2}{S_1 \sqcap S_2}}$ (with $\top$ absorbing on
either side, kept exactly), \emph{given} the STypes-level containment;
the \texttt{None} result inherits its status from
\texttt{STypes::inter}, which is exact (Section~\ref{sec:inter-stypes}),
so the clause-level bottom is semantic.
\end{lemma}
\begin{proof}
$(\bigcap_{V_1} \cap\, \sem{S_1}) \cap (\bigcap_{V_2} \cap\, \sem{S_2})
= \bigcap_{V_1 \cup V_2} \cap\, (\sem{S_1} \cap \sem{S_2})$, then the
STypes containment. The $\top$ arms are the identities
$\Val \cap X = X$.
\impl{TTypeInter.sk / Clause::inter}
\end{proof}

\begin{lemma}[Clause Product and Collision Re-Union]
\label{lem:clauses-inter}
$\sem{C_1} \cap \sem{C_2} \subseteq
\sem{\Clauses\mathord{::}\inter(C_1,C_2)}$ given
Lemma~\ref{lem:clause-inter}'s hypothesis, and the constructor invariant
is preserved.
\end{lemma}
\begin{proof}
$\sem{C_1} \cap \sem{C_2} = \bigcup_{c_1, c_2} \sem{c_1} \cap \sem{c_2}$
(distributivity), and each pair lands in the result: directly, or merged
by \texttt{Clause::union} when two pairs produce the same var-set key —
exact by Lemma~\ref{lem:clause-union}. Dropped (\texttt{None}) pairs
contribute nothing by Lemma~\ref{lem:clause-inter}. Invariant: the
result is built without the top-collapsing constructor, but a top clause
$\clause{\emptyset}{\top}$ can arise only from a pair of top clauses,
and a top clause is alone in its operand (invariant), so both operands
were top and the result is the lone top clause. An empty product yields
empty $\Clauses$, i.e.\ $\bot$ at the $\TType$ level — semantic iff each
pair's bottom was.
\impl{TTypeInter.sk / Clauses::inter}
\end{proof}

\subsection{The instance layer: TVar constraints, TInter, speculative
returns}
\label{sec:inter-inst}

\texttt{inter\_} has five early arms before the clause product; we
audit each.

\paragraph{(i) Two unbound bare variables.}
Returns the unexpanded intersection
$\mathsf{TInter}\{\alpha, \beta\}$, whose interpretation \emph{is}
$\sem{\alpha} \cap \sem{\beta}$ (the class comment's
$\sqcap$-interpretation; its \texttt{expand} recomputes the intersection
under the then-current state, covered by \textbf{L-EXP}). Containment is
definitional; never $\bot$.
\impl{TTypeInter.sk / TInter (interpretation comment), TType::unexpandedInter}

\paragraph{(ii)/(iii) One bare variable, unbound or open-for-field.}
Returns $\alpha$ with the other operand attached as a constraint:
\texttt{TVar::setConstraint} is a \emph{functional} update — \texttt{TVar}
is frozen, the method returns a new \texttt{TVar} value, and constraints
are occurrence-level, not shared state (registry row SETCONSTRAINT-FUN,
deep dive~E) — so no other occurrence of $\alpha$ changes meaning. For
the returned occurrence, the main paper's canonical TVAR-SEM definition
states
$\sem{\alpha^{\mathit{cs}}}_\rho = \sem{\alpha^{\emptyset}}_\rho
\cap \bigcap_{c \in \mathit{cs}} \sem{c}_\rho$, matching
\texttt{Subst::get}, which intersects the binding's expansion with the
occurrence's constraints at every read. Then
$\sem{\alpha^{\mathit{cs}}} \cap \sem{t} =
\sem{\alpha^{\mathit{cs} \cup \{t\}}}$ \emph{exactly}, at every state —
provided the get-time intersection itself over-approximates, which is
L-INTER/L-INTER-BOT again, applied at read time; the circularity is benign (an
inductive appeal at the strictly later read). Never $\bot$.
\impl{TTypeInter.sk / inter\_ (setConstraint arms); TType.sk /
TVar::setConstraint; TTypeSubst.sk / Subst::get (constraint loop)}

\paragraph{(iv) In-flight cycle.}
If the pair is already being intersected (\texttt{tenv.seenInter}, an
in-flight \emph{set}, distinct from the memo cache), return the
unexpanded intersection — same justification as (i); never $\bot$.

\paragraph{(v) Speculative boxed-primitive returns.}
When one side expands to a single primitive kind and the other to a
single object type, the code checks
$\mathsf{isSub}(\mathsf{Box}(k), \mathsf{object})$ and on success
returns the \emph{primitive operand unchanged} (e.g.\
\texttt{number \& \{toFixed: ...\} = number}). Returning either operand
always over-approximates the intersection, so these arms are sound
regardless of the subtype verdict, and they never produce $\bot$. (They
do evolve the state through \texttt{isSub}'s growth; the code comments
note only concrete types are involved, and failure paths restore — the
clone/restore discipline of the main paper.)
\impl{TTypeInter.sk / inter\_ (primKind/getObject arms)}

\paragraph{Expansion drops nothing.}
The clause product runs on
\texttt{expandToClauses(subst, tenv, Covariant())}, which discards the
first component of \texttt{expandToVectors}. We verified that component
is populated \emph{only} when \texttt{expandTVars = None()} (variables
pushed back unexpanded); with a variance supplied, every instance —
\texttt{TSInst} (via the symbolically-cached expansion),
\texttt{TInter}, \texttt{TArrayLit}, \texttt{TVar} — is rewritten into
clauses, so the discarded vector is empty and no members are lost.
A top expansion short-circuits to $\top$, exact for intersection's
identity.
\impl{TType.sk / expandToVectors, expandToClauses}

\subsection{The bottom inventory}
\label{sec:bottom-inventory}

Every site where intersection can produce $\bot$ (\texttt{None} or an
empty structure), with its verdict. ``Semantic'' means the empty result
implies empty semantic intersection; ``derived'' means the verdict is
inherited from the sites it aggregates.

\begin{center}
\renewcommand{\arraystretch}{1.25}
\small
\begin{tabular}{@{}p{4.9cm}p{5.6cm}p{3.4cm}@{}}
\toprule
\textbf{Site} & \textbf{When $\bot$} & \textbf{Nature of $\bot$} \\
\midrule
\texttt{PrimLiterals::inter} & spelling sets disjoint &
  semantic iff L-LIT(ii); exact for decoded strings and
  within-precision numbers; spurious for beyond-precision num
  spellings (Finding~\ref{find:spelling}) \\
\texttt{TPrimValue::inter} & kind mismatch & semantic (A1); unreachable
  from \texttt{STypes::inter} \\
\texttt{STypes::inter}, no shared tag survives & all per-tag results
  empty or unmatched & semantic (A1), derived \\
\texttt{STypes::inter}, empty shared tag & a shared tag's per-tag
  intersection is empty & semantic (A1); dropped, not zeroed \\
\texttt{TFields::inter}, brand & $\mathsf{Some}(\mathsf{true})$ vs
  $\mathsf{Some}(\mathsf{false})$ & semantic \\
\texttt{interCallability}, array-branded conflict &
  $\mathsf{CAbsent} \sqcap \mathsf{CPresent}/\mathsf{CFunction}$ under
  $\mathsf{isArray} = \mathsf{Some(true)}$ & semantic, conditional on
  the array brand's membership invariant (F-BRAND) \\
\texttt{interCallability}, signature arm & \texttt{TFuns::inter} fails
  both \texttt{isSub} directions
  ($\mathsf{CPresent} \sqcap \mathsf{CPresent}$) & heuristic
  (Finding~\ref{find:tfun}) \\
\texttt{TFields::interIsArray} & brand conflict at \textsc{R-IsArray} &
  semantic \\
\texttt{TObject::inter} & every branch pair \texttt{None} & derived \\
\texttt{Clause::inter} / \texttt{Clauses::inter} & STypes bottom /
  empty product & derived \\
\texttt{TInter::map} & a mapped component is $\bot$ & derived (client
  utility, not on the \texttt{inter} path) \\
\texttt{inter\_} short-circuit & an operand is $\bot$ & semantic \\
\bottomrule
\end{tabular}
\end{center}

\noindent
Notably, \emph{fields and indexes never produce} $\bot$: two named
fields with disjoint gets intersect to a field whose get is $\bot$ —
a branch that no object with that key present can inhabit, but the
branch (and, if the field is optional, its members lacking the key)
survives. The implementation's own \texttt{overlaps} comment documents
this deliberately (``a shared field with non-overlapping types'' is
listed as a TODO for \emph{warning} precision, not as a bottom). So the
only unsound bottoms are the signature arm
(Finding~\ref{find:tfun}) and the beyond-precision range of the
literal stratum (Finding~\ref{find:spelling}) — but each sits on a hot
refinement path.

%==========================================================================
\section{The Caches: seenUnion and seenInter}
\label{sec:caches}

Both caches live in the substitution state, keyed by the operand pair,
while the substitution evolves between write and reuse. The question is
whether a cached result can become an \emph{under}-approximation after
evolution (unsoundness), as opposed to merely stale precision.

\begin{theorem}[seenUnion Reuse is Unconditionally Sound]
\label{thm:seenunion}
A \texttt{seenUnion} entry reused at any later state $\rho'$ still
satisfies the L-UNION containment at $\rho'$.
\end{theorem}
\begin{proof}
By the proof of Theorem~\ref{thm:union}, \texttt{union\_} never reads
$\rho$: it is a pure function of the operands' syntax, and the
containment proved is pointwise in $\rho$ — the result type's members
are (unions of) the very same syntactic components, so
$\sem{t_1}_{\rho'} \cup \sem{t_2}_{\rho'} \subseteq
\sem{\union(t_1,t_2)}_{\rho'}$ holds for every $\rho'$ with one proof.
The cached value equals the recomputation verbatim, so reuse is the
identity. (This is the code comment's ``TVars pass through unchanged''
argument, made precise: state-independence of the computation plus
state-uniformity of the theorem.) \texttt{restore} additionally rolls
the cache back together with the bindings, which is needed only for
precision hygiene, not soundness.
\impl{TTypeUnion.sk / TType::union (cache and comment);
TTypeSubst.sk / SubstState.seenUnion, restore}
\end{proof}

\subsection{seenInter: the sealing argument and its freeze exception}

\texttt{seenInter} is different in kind: \texttt{inter\_} both reads and
evolves the state. The implementation's justification (the comment on
\texttt{TType::inter}) is a \emph{sealing} argument:
\begin{enumerate}
\item Bare-variable operands bypass the cache entirely (the
  \texttt{hasTVar} guard) — exactly the arms of
  Section~\ref{sec:inter-inst} that touch bindings without sealing.
\item Every variable a cached computation \emph{did} read was read
  through \texttt{expandToClauses(\ldots, Covariant())}, i.e.\
  \texttt{Subst::get} covariantly, which seals it
  ($\mathsf{SVOpen} \to \mathsf{SVClosed}$; fresh $\to$
  $\mathsf{SVClosed}(\top)$); a sealed binding never changes again (main
  paper, Lemma ``Determined Candidates Are Never Widened''(iii), which
  holds for all identifiers), so the expansion the cache baked in equals
  the expansion any later state would produce.
\item Named-type expansions consulted along the way
  (\texttt{seenExpansion}) are symbolic in their type arguments — the
  cached body mentions the argument \emph{variables}, whose bindings are
  re-resolved at every use — so they cannot freeze a binding either
  (the argument documented at \texttt{SubstState.seenExpansion}).
\item \texttt{restore} rolls back \texttt{seenInter} together with the
  bindings, so failure-path exploration cannot leak entries into a state
  they were not computed against; and a \texttt{clone} shares the maps
  by value (immutable), so a clone's writes never reach the parent.
\end{enumerate}
This argument is correct as far as it goes, but it has one exception,
which we found by instruction-level audit of \texttt{Subst::get}: under
\texttt{freezeCandidates}, a covariant read of an $\mathsf{SVOpen}$
\emph{candidate} returns the value \emph{without sealing}. An
\texttt{inter} executed inside a constraint-admissibility check (the
freeze window of \texttt{growTVars} — reachable, since admissibility
calls \texttt{isSub}, whose validation intersects variable constraints
through \texttt{Subst::get}) can therefore cache a result derived from a
binding that remains growable after the freeze is popped; the freeze
flag is deliberately outside \texttt{SubstState}, so no restore removes
the entry. If the candidate later grows (a contravariant union enlarging
its denotation), the cached intersection under-approximates the new
$\sem{t_1} \cap \sem{t_2}$.

\begin{obligation}[C-FREEZE-CACHE: seenInter Entries Written Under
Freeze]
\label{obl:freeze-cache}
The mechanism is \emph{demonstrably active}: \texttt{TType::inter} is
reached inside freeze windows with unsealed candidates (via
\texttt{Subst::get}'s constraint-application loop), \texttt{seenInter}
entries are written there, and an end-to-end order-dependent unsound
acceptance exercising this path exists — shared with the sibling
\texttt{subst.seenSub} memo (deep dive E's MR-1, \textsc{violated};
issue draft \path{issues/freeze-window-stale-memo-caches.md}). In
particular, two otherwise-tempting deflections are false: admissibility
checks do \emph{not} save/restore around themselves on success paths,
and freeze-window operand pairs do \emph{not} always carry a bare
variable. What remains open is \emph{attribution} only: no witness
isolates \texttt{seenInter}'s contribution from \texttt{seenSub}'s
(every construction that reused a stale \texttt{seenInter} entry also
reused a stale \texttt{seenSub} entry). Discharge the isolated
statement, or apply the enforcing fix, which covers both memos: key
\texttt{seenInter} by substitution version (as \texttt{tenv.seenSub}
already is, after \#446 — \texttt{subst.seenSub} itself is \emph{not}
version-keyed and shares this exposure), or skip cache \emph{writes}
while frozen.
\impl{TTypeSubst.sk / Subst::get (freeze arm), Subst.freezeCandidates;
TTypeInter.sk / TType::inter (hasTVar guard, re-entrant arm)}
\end{obligation}

\begin{lemma}[Constraint Attachment is Exact under TVAR-SEM]
\label{lem:tvar-sem}
By the main paper's canonical variable denotation (registry row
\textbf{TVAR-SEM}),
$\sem{\alpha^{\mathit{cs}}}_\rho = \sem{\alpha^{\emptyset}}_\rho \cap
\bigcap_{c}\sem{c}_\rho$, with
$\sem{\alpha^{\emptyset}}_\rho=\emptyset$ when $\alpha$ is unbound.
Therefore Section~\ref{sec:inter-inst}(ii)'s constraint-attachment result
is exact.
\end{lemma}

Two further hygiene points, verified benign: the re-entrant arm of
\texttt{TType::inter} (a recursive call cached the pair first) returns
the \emph{already-cached} result rather than the just-computed one, for
consistency with intermediate consumers — soundness is unaffected since
both are results of \texttt{inter\_} on the same pair; and the
\texttt{ty1 == ty2} fast path returns \texttt{ty1}, exact for
idempotent operands.

%==========================================================================
\section{Empirical Witnesses}
\label{sec:findings}

All probes were run against the pinned binary
(\texttt{target/host/dev/skjs}) via
\texttt{make check-file-in-docker}. ``Accepted'' means zero diagnostics.

\subsection{Probe 2: signature bottom (Finding~\ref{find:tfun})}
\texttt{if (k === g)} on incomparable signature types narrows \texttt{k}
to \texttt{never} in a branch a single accepted program takes at
runtime; the end-to-end witness and the tsc comparison are issue draft
\path{issues/tfun-inter-incomparable-not-disjoint.md}.

\subsection{Probe 3: literal spellings (Finding~\ref{find:spelling})}
\begin{finding}[Spelling-Sensitive Disjointness]
\label{find:spelling}
Beyond-precision numeric spellings break \textbf{L-INTER-BOT}:
reading: spellings that alias one double are kept distinct (the
num-kind \textbf{L-LIT}(ii) violation, issue draft
\path{issues/beyond-precision-numeric-literals.md}, which carries the
executed witness: two $2^{53}$ spellings make an accepted
\texttt{===}-true branch live at runtime with a spurious \texttt{never},
and a number flows into \texttt{string}), so their lexeme intersection
is empty while their value sets coincide — the failure propagating
through Lemma~\ref{lem:lit-inter} into \texttt{PrimLiterals::inter}'s
bottom and thence into \texttt{overlaps} and the branch logic.
\emph{Within-precision} numeric spellings do not trigger it: the parser
normalizes them, so \texttt{0.0} and \texttt{0} share one lexeme.
\emph{String} spellings do not trigger it either: escape sequences are
decoded at parse time, so distinct stored lexemes denote distinct
strings and the spelling bottom is semantic
(Lemma~\ref{lem:lit-inter}).
The comparison \emph{lint} knows the runtime truth —
\texttt{checkComparisonOverlap} is silenced for imprecise pairs by
\texttt{hasImpreciseNumberLiteral}
(\texttt{TypeCheckExpr.sk:1656-1657, 1704}), its comment citing this
very pair — but the narrowing and intersection path carries no
precision guard, so the spurious $\bot$ is still installed.
\end{finding}

\subsection{Probe 4: type-level intersections}
\begin{verbatim}
type C = ("a" | 0) & ("a" | 1);
function p6(): void { const c: C = "a"; }        // ACCEPTED: C = "a"
type F = ((s: string) => void) & ((n: number) => void);
function p8(f: F): void { const z: never = f; }  // ACCEPTED: F = never
\end{verbatim}
The STypes assembly is exact: \texttt{C = "a"} (the empty \texttt{num}
tag is dropped, not zeroed), matching tsc. The call/new stratum still
collapses: \texttt{F} is \texttt{never} where tsc reads a callable
overload — \texttt{p8}'s acceptance is the incompleteness dual of
Finding~\ref{find:tfun}.

\subsection{Failed proof attempts, recorded}
For completeness, the containments we attempted and could not close,
each of which became a Finding or Obligation rather than a theorem:
(1) call/new-level $\bot$-implies-disjoint — refuted
(Finding~\ref{find:tfun});
(2) literal-level $\bot$-implies-disjoint without the beyond-precision
num-kind L-LIT(ii) — refuted for that range
(Finding~\ref{find:spelling});
(3) unconditional \texttt{seenInter} reuse — the sealing argument's
freeze gap resisted proof and is Obligation~\ref{obl:freeze-cache}.
We found \emph{no} defect in: any union stratum, the field/index/brand
strata of intersection, the instance-layer arms, the clause product, or
the in-flight cycle handling.

%==========================================================================
\section{Status and Ledger Edits}
\label{sec:status}

\subsection{Status by stratum}

\begingroup
\renewcommand{\arraystretch}{1.25}
\small
\begin{longtable}{@{}p{4.6cm}p{9.4cm}@{}}
\toprule
\textbf{Stratum} & \textbf{Result} \\
\midrule
\endhead
Union: literals, prims, units & Lemmas~\ref{lem:lit-union},
  \ref{lem:prim-union} (exact) \\
Union: object branch sets & Lemma~\ref{lem:obj-union} (exact) \\
Union: STypes / Clause / Clauses / TInsts &
  Lemmas~\ref{lem:stypes-union}--\ref{lem:insts-union},
  Theorem~\ref{thm:union} (exact) \\
Union: per-read merge clients (\texttt{TField::union},
  \texttt{modeOf}) & Lemma~\ref{lem:tfields-union} (over-approx.) \\
\texttt{seenUnion} cache & Theorem~\ref{thm:seenunion} (unconditional) \\
\midrule
Inter: units, tag dropping, no-shared-tag $\bot$ &
  Lemma~\ref{lem:prim-inter}, A1 \\
Inter: literal sets & Lemma~\ref{lem:lit-inter}, conditional on
  L-LIT(ii); exact for decoded strings and within-precision numbers;
  beyond-precision num spellings spurious
  (Finding~\ref{find:spelling}) \\
Inter: STypes assembly & exact — empty shared tags dropped, $\bot$ iff all
  drop \\
Inter: named fields, indexes, brand, access &
  Section~\ref{sec:inter-object} (brand $\bot$ semantic; given
  sub-level L-INTER/L-INTER-BOT) \\
Inter: callability meet & sound outside the signature arm (array-brand
  $\bot$ semantic modulo F-BRAND); signature-failure branch drop
  (Finding~\ref{find:tfun}, pinned probe) \\
Inter: Clause/Clauses product & Lemmas~\ref{lem:clause-inter},
  \ref{lem:clauses-inter} (derived) \\
Inter: instance layer (constraints, TInter, speculative returns) &
  Section~\ref{sec:inter-inst}; never $\bot$, via TVAR-SEM
  (Lemma~\ref{lem:tvar-sem}) \\
\texttt{seenInter} cache & sealing argument + freeze gap
  (Obligation~\ref{obl:freeze-cache}) \\
Cloned-subst reuse of results & clone-stability, via the fork-state
  form I$_f$ of the containments audited here \\
\bottomrule
\end{longtable}
\endgroup

\subsection{Stratum outcomes}

\paragraph{L-UNION.} \textsc{proved}
(this document, Theorem~\ref{thm:union} with
Lemmas~\ref{lem:lit-union}--\ref{lem:tfields-union} and
Theorem~\ref{thm:seenunion} for the cache). Note field: ``exact for
membership semantics at every stratum on the $\union(t_1,t_2)$ path;
the per-read merge clients (\texttt{TField::union}, the
\texttt{modeOf} callability join) are over-approximating only.
Signature bodies built by \texttt{TFun::union} — including its
\texttt{this}-parameter intersection — are outside membership
semantics and remain part of the subtyping surface.''

\paragraph{L-INTER and L-INTER-BOT.} The two cases are graded separately:
\begin{itemize}
\item \textbf{L-INTER} is \textsc{proved}, conditional on
  \textbf{L-INTER-BOT} for recursive calls, whenever a fresh structural
  intersection returns non-bottom: units, tags, literal sets, fields,
  indexes, brand, callability, clause product, STypes assembly, and the
  instance layer (the last using \textbf{TVAR-SEM}). This recursive premise
  is necessary: a non-bottom outer object may contain a field result
  produced by a nested intersection that spuriously returned $\bot$.
\item \textbf{L-INTER-BOT} is \textsc{proved} for units, tags, decoded
  strings, within-precision numbers, fields, indexes, brand (modulo
  F-BRAND), clause product, STypes assembly, and the instance layer. It is
  \textsc{violated} for beyond-precision numeric spellings and for the
  callability signature arm (Findings~\ref{find:spelling}
  and~\ref{find:tfun}).
\item Cached \texttt{seenInter} reuse remains \textsc{open} for both
  guarantees where inference growth can make a retained result stale
  (C-FREEZE-CACHE, Obligation~\ref{obl:freeze-cache}). Consequently the
  registry grades L-INTER \textsc{open}, with the fresh structural case
  proved; L-INTER-BOT is already \textsc{violated} by direct witnesses.
\end{itemize}
Every theorem of the main paper that needs an inhabited intersection to
survive consumes both rows. The
branch-drop steps of \textsc{R-Eq}, \textsc{R-EqField},
\textsc{R-Inst}, \textsc{R-IsArray}, and the equality-narrowing
intersection itself are affected by the violated L-INTER-BOT ranges
exactly where its operands reach them; the \textsc{R-IsArray} drops
are unaffected (their bottoms are the brand and no-shared-tag cases,
both semantic — except when the intersected array branch pair hits the
signature arm, which \texttt{ReadonlyArray} does not). As with the
main paper's other \textsc{violated} entries, the violations are local
to leaf strata: both are code fixes after which the structural
L-INTER-BOT cases are all \textsc{proved} (modulo F-BRAND's brand
invariant at its array arm), and only the \textsc{open} obligation
C-FREEZE-CACHE remains here. Each violated
stratum carries a pinned end-to-end repro with its issue draft under
\path{issues/repros/}, asserted by the registry tooling.

\paragraph{Registry.}
The authoritative, current statuses are maintained in the registry
(\texttt{docs/ledger.yaml}), rendered as the main paper's Assumptions and
Proof Obligations tables; this paper is cited there as evidence. The
registry rows it supports:

% GENERATED by docs/tools/docs.py from docs/ledger.yaml — do not edit.

\begin{center}
\renewcommand{\arraystretch}{1.25}
\small
\begin{tabular}{@{}>{\raggedright\arraybackslash}p{2.4cm}>{\raggedright\arraybackslash}p{8.2cm}>{\raggedright\arraybackslash}p{2.6cm}@{}}
\toprule
\textbf{Row} & \textbf{Statement} & \textbf{Status} \\
\midrule
L-UNION & \texttt{union} over-approximates set union (all levels). & \textsc{proved} \\
L-INTER & When \texttt{inter} returns a non-$\bot$ result, that result contains the semantic intersection of its operands, provided every recursive call that returned $\bot$ had disjoint operands (L-INTER-BOT). & \textsc{open} \\
L-INTER-BOT & A $\bot$ result from \texttt{inter} implies semantic disjointness. & \textsc{violated} \\
C-FREEZE-CACHE & A stale \texttt{seenInter} entry can cause an unsound acceptance after candidate growth resumes without also reusing a stale \texttt{subst.seenSub} entry. & \textsc{open} \\
\bottomrule
\end{tabular}
\end{center}


\paragraph{Imports (review contract).}
The named statements this paper assumes from elsewhere. A reviewer treats
each as an axiom, checking citations against the registry's canonical
statement; no other paper's internals are in scope:

% GENERATED by docs/tools/docs.py from docs/ledger.yaml — do not edit.

\begin{center}
\renewcommand{\arraystretch}{1.25}
\small
\begin{tabular}{@{}>{\raggedright\arraybackslash}p{2.6cm}>{\raggedright\arraybackslash}p{7.4cm}>{\raggedright\arraybackslash}p{1.2cm}>{\raggedright\arraybackslash}p{2.2cm}@{}}
\toprule
\textbf{Row} & \textbf{Canonical statement} & \textbf{Owner} & \textbf{Status} \\
\midrule
L-EXP & Expansion preserves or enlarges semantics (modes \texttt{None}/Covariant); residual instances keep theirs. & C & \textsc{proved} \\
L-LIT & Literal spellings: (i) $\delta$ is a function; (ii) injectivity modulo $\approx$ after normalization; (iii) falsiness-exactness of the lexeme test. & G & \textsc{violated} \\
MR-1 & No \texttt{subst.seenSub} memo recorded under \texttt{freezeCandidates} may outlive growth of an open candidate readable from its left side. & E & \textsc{violated} \\
TVAR-SEM & A type-variable occurrence denotes its substitution envelope intersected with its occurrence constraints; an unbound envelope is empty. & main & \textsc{proved} \\
F-BRAND & Subtyping never compares the array brand, though refinement narrows on it. & D & \textsc{violated} \\
SETCONSTRAINT-FUN & \texttt{TVar::setConstraint} is a functional update on a frozen value: constraints are occurrence-level, not shared state. & E & \textsc{proved} \\
\bottomrule
\end{tabular}
\end{center}


%==========================================================================
\section{Implementation Anchor Index}
\label{sec:anchors}

Pinned to \texttt{\skjspinhash}.

\begin{center}
\renewcommand{\arraystretch}{1.2}
\small
\begin{tabular}{@{}p{4.1cm}>{\raggedright\arraybackslash}p{10.3cm}@{}}
\toprule
\textbf{This document} & \textbf{Implementation} \\
\midrule
union top level, cache & \texttt{TTypeUnion.sk / TType::union,
  union\_} \\
literal-set ops & \texttt{PrimLiterals.sk / union, inter, diff};
  \texttt{JSStringEscape.sk / decodeJSStringLiteral} (decode at parse) \\
primitive merges & \texttt{TTypeUnion.sk / TPrimValue::union};
  \texttt{TTypeInter.sk / TPrimValue::inter} \\
STypes assembly & \texttt{TTypeUnion.sk / STypes::union};
  \texttt{TTypeInter.sk / STypes::inter} \\
map combinators & \texttt{InternalKeyedArrayMap.sk / mapUnion,
  mapInterOption, mergeOption, mergeOpt} \\
object branch sets & \texttt{TTypeUnion.sk / TObject::union};
  \texttt{TTypeInter.sk / TObject::inter} \\
branch pair & \texttt{TTypeInter.sk / TFields::inter,
  interNamedFields, TIndexes::inter} \\
per-read merges & \texttt{TType.sk / TObject::maybeGetType,
  resolveLiteralAccess, maybeGetSymbol, modeOf, callMode, newMode,
  maybeGetCall/New, maybeGetAccessibleIndexForKeyKind} \\
fields & \texttt{TTypeInter.sk / TField::inter, interIndex};
  \texttt{TTypeUnion.sk / TField::union, unionIndex} \\
access lattice & \texttt{TType.sk / FieldAccess::combine} \\
callability & \texttt{TType.sk / Callability,
  TFields::createFromSortedIterator (invariants)};
  \texttt{TTypeInter.sk / TFields::interCallability, TFun::inter,
  TFuns::inter}; \texttt{TTypeUnion.sk / TFun::union (incl.\
  thisParam), TFuns::union} \\
brand & \texttt{TTypeInter.sk / TFields::inter (isArray),
  interIsArray} \\
clause strata & \texttt{TTypeUnion.sk} / \texttt{TTypeInter.sk} /
  \texttt{Clause::*}, \texttt{Clauses::*}; \texttt{TType.sk /
  Clauses::create} \\
instance layer & \texttt{TTypeInter.sk / TInter, inter\_ (var arms)};
  \texttt{TType.sk / TVar::setConstraint} \\
expansion & \texttt{TType.sk / expandToVectors, expandToClauses} \\
caches & \texttt{TTypeSubst.sk / SubstState.seenUnion, seenInter,
  seenExpansion, restore; Subst::clone, get, freezeCandidates} \\
overlap warning & \texttt{TTypeInter.sk / SType::overlaps} \\
refinement consumers & \texttt{TypeCheckRefine.sk / evalVarEquality,
  evalFieldEqualityType, narrowField, evalIsArray, evalInstanceof} \\
\bottomrule
\end{tabular}
\end{center}

\end{document}
